% ============================================================
%  翻译工作区 · 通用排版导言区
%  引擎：XeLaTeX（必须）
%  项目专用改动请写在 project.tex 或正文中，尽量不改本文件
% ============================================================

% ---------- 宏包 ----------
\usepackage{geometry}
\usepackage{fancyhdr}
\usepackage{enumitem}
\usepackage{booktabs}
\usepackage{longtable}
\usepackage{array}
\usepackage{tabularx}
\usepackage{multirow}
\usepackage{graphicx}
\usepackage{xcolor}
\usepackage{float}
\usepackage{setspace}
\usepackage{ragged2e}
\usepackage{caption}
\usepackage{tocloft}
\usepackage{amsmath,amssymb}
\usepackage{microtype}

% ---------- 页面 ----------
\geometry{
  paperwidth=170mm,
  paperheight=240mm,
  top=25mm,
  bottom=25mm,
  inner=22mm,
  outer=18mm,
  headheight=14pt,
  headsep=6mm,
  footskip=12mm
}

% ---------- 字体 ----------
% 西文：Palatino（正文），Consolas（等宽）
\setmainfont{Palatino Linotype}[
  BoldFont={Palatino Linotype Bold},
  ItalicFont={Palatino Linotype Italic},
  BoldItalicFont={Palatino Linotype Bold Italic}
]
\setsansfont{Segoe UI}
\setmonofont{Consolas}[Scale=MatchLowercase]

% 中文：Noto Serif SC（正文）、Noto Sans SC（标题/无衬线）、楷体、仿宋
% 注：系统安装的 Noto SC 为可变字体（VF），xdvipdfmx 无法嵌入，
% 故正文用宋体、标题用黑体；两者均为 Windows 自带、可稳定嵌入。
% 如安装了静态版思源宋体/黑体，可改为 Source Han Serif SC / Source Han Sans SC。
\setCJKmainfont{SimSun}[AutoFakeBold=1.5]
\setCJKsansfont{SimHei}
\setCJKmonofont{FangSong}
\setCJKfamilyfont{zhsong}{SimSun}[AutoFakeBold=1.5]
\setCJKfamilyfont{zhhei}{SimHei}
\setCJKfamilyfont{zhkai}{KaiTi}
\setCJKfamilyfont{zhfs}{FangSong}
\providecommand{\songti}{}
\providecommand{\heiti}{}
\providecommand{\kaishu}{}
\providecommand{\fangsong}{}
\renewcommand{\songti}{\CJKfamily{zhsong}}
\renewcommand{\heiti}{\CJKfamily{zhhei}}
\renewcommand{\kaishu}{\CJKfamily{zhkai}}
\renewcommand{\fangsong}{\CJKfamily{zhfs}}

% ---------- 行距与段落 ----------
\linespread{1.45}
\setlength{\parindent}{2em}
\setlength{\parskip}{0pt}
\setlength{\emergencystretch}{2em}

% ---------- 章节标题 ----------
\ctexset{
  chapter = {
    format = \centering,
    nameformat = \sffamily\large\bfseries,
    numberformat = \sffamily\large\bfseries,
    titleformat = \sffamily\LARGE\bfseries,
    name = {第,章},
    number = \chinese{chapter},
    beforeskip = 20pt,
    afterskip = 40pt
  },
  section = {
    format = \sffamily\large\bfseries,
    beforeskip = 1.4ex plus .3ex minus .2ex,
    afterskip = .8ex plus .2ex
  },
  subsection = {
    format = \sffamily\normalsize\bfseries,
    beforeskip = 1.2ex plus .2ex minus .1ex,
    afterskip = .6ex plus .1ex
  }
}
% 章首不显示页眉线，章标题页使用 plain 样式
\renewcommand{\chaptermark}[1]{\markboth{#1}{#1}}

% 空白页不显示页眉页脚（book 类默认会给空白页加页眉）
\makeatletter
\renewcommand{\cleardoublepage}{%
  \clearpage
  \if@twoside
    \ifodd\c@page\else
      \hbox{}\thispagestyle{empty}\newpage
      \if@twocolumn\hbox{}\newpage\fi
    \fi
  \fi}
\makeatother

% 目录、插图目录标题居中
\renewcommand{\contentsname}{\hfill 目\quad 录\hfill}
\renewcommand{\listfigurename}{\hfill 插图目录\hfill}

% ---------- 页眉页脚 ----------
\pagestyle{fancy}
\fancyhf{}
\fancyhead[LE]{\small\itshape\booktitle}
\fancyhead[RO]{\small\itshape\leftmark}
\fancyfoot[C]{\small\thepage}
\renewcommand{\headrulewidth}{0.4pt}
\renewcommand{\footrulewidth}{0pt}
\fancypagestyle{plain}{%
  \fancyhf{}%
  \fancyfoot[C]{\small\thepage}%
  \renewcommand{\headrulewidth}{0pt}%
}

% ---------- 脚注 ----------
\usepackage[bottom,hang,flushmargin]{footmisc}
\renewcommand{\footnotelayout}{\small}
% 原书脚注每章从 1 开始，与原文编号对应
\counterwithin*{footnote}{chapter}
\setlength{\skip\footins}{1.1\baselineskip}

% ---------- 引文 ----------
% 短引文：随文引号；长引文：quote 环境（自动缩进 + 小字号）
\renewenvironment{quote}
  {\list{}{\leftmargin=2em\rightmargin=2em}\item[]\small\kaishu}
  {\endlist}

% 章首题词（题词正文译出，出处保留原文）
% 用法：\begin{chapterquote}{——作者，\textit{书名}}
%          ……题词……
%       \end{chapterquote}
\newenvironment{chapterquote}[1]
  {\def\chapterquotesource{#1}\begin{quote}\kaishu\small}
  {\par\vspace{0.8em}\hfill\upshape\small\chapterquotesource\end{quote}}

% ---------- 表格 ----------
\renewcommand{\arraystretch}{1.25}
\setlength{\tabcolsep}{6pt}
\captionsetup[table]{position=top}

% ---------- 图注 ----------
\captionsetup{
  font=small,
  labelfont=bf,
  labelsep=quad,
  skip=6pt
}
\graphicspath{{figures/}{source/images/}}

% ---------- 列表 ----------
\setlist[itemize]{leftmargin=2em,itemsep=0.2ex,topsep=0.4ex}
\setlist[enumerate]{leftmargin=2em,itemsep=0.2ex,topsep=0.4ex}

% ---------- 译者工具宏 ----------
% 术语首次出现：中文（原文）
% 用法：\term{述行性}{performativity}
\newcommand{\term}[2]{#1（\textit{#2}）}

% 仅保留原文（限定为西文斜体）
\newcommand{\orig}[1]{\textit{#1}}

% 译者注：独立编号（每章重新计数），不占用原书脚注编号。
% 页脚显示“译者注1：……”，正文标记为“〔译注1〕”
% 用法：\zhu{……}
\newcounter{zhu}[chapter]
\newcommand{\zhu}[1]{%
  \stepcounter{zhu}%
  \textsuperscript{〔译注\thezhu〕}%
  \begingroup
  \renewcommand{\thefootnote}{}%
  \footnotetext{\textbf{译者注\thezhu：}#1}%
  \endgroup}


% 独立参考文献/原注中的书名等，避免过度排版
\newcommand{\bookorig}[1]{\textit{#1}}

% 待办标记：编译可见，终检前必须清零
\newcommand{\todo}[1]{\textcolor{red}{【TODO：#1】}}

% 译者按语（块级，带底纹）
\usepackage[breakable]{tcolorbox}
\newtcolorbox{transnote}[1][译者注]{
  breakable,
  colback=black!3,
  colframe=black!35,
  boxrule=0.4pt,
  arc=1.5pt,
  left=8pt,
  right=8pt,
  top=6pt,
  bottom=6pt,
  fonttitle=\small\sffamily\bfseries,
  title={#1},
  before skip=1.2ex,
  after skip=1.2ex
}

% 图注来源行
% 用法：\figsource{图片来源：某某出版社}
\newcommand{\figsource}[1]{\par\vspace{0.3em}{\footnotesize #1\par}}

% ---------- 目录 ----------
\renewcommand{\cftchapfont}{\sffamily\bfseries}
\renewcommand{\cftchappagefont}{\sffamily}
\setlength{\cftbeforechapskip}{0.6em}
\setlength{\cftsecindent}{2em}

% ---------- 超链接与 PDF 元数据 ----------
\usepackage[
  colorlinks=true,
  linkcolor=black,
  citecolor=black,
  urlcolor=blue!45!black,
  bookmarksnumbered=true,
  bookmarksopen=true,
  pdfencoding=auto
]{hyperref}
\hypersetup{
  pdftitle={\booktitle},
  pdfauthor={\bookauthor},
  pdfsubject={\booktitleoriginal},
  pdfkeywords={\bookkeywords},
  pdfcreator={XeLaTeX}
}

% ---------- 其他 ----------
